\documentclass[12pt]{article}
\usepackage[utf8]{inputenc}
\usepackage[T1]{fontenc}
\usepackage[margin=1in]{geometry}
\usepackage{setspace}
\usepackage{xcolor}
\usepackage[colorlinks=true, urlcolor=blue!60!black, linkcolor=blue!60!black]{hyperref}
\setlength{\parindent}{0em}
\setlength{\parskip}{0.6em}
\renewcommand{\baselinestretch}{1.0}
\begin{document}
\begin{center}
{\LARGE \textit{Shelley v. Kraemer}, 334 U.S. 1 (1948)\par}
\vspace{0.5em}
{\large Case Brief --- Two-Page Edition\par}
\end{center}

\section{Memory Jogger}
Racially restrictive covenants are private, but state-court enforcement of them is state action that denies equal protection (\href{https://supreme.justia.com/cases/federal/us/334/1/}{Justia}).

\section{Detailed Case Facts}
Recorded covenants in St. Louis and Detroit barred occupancy by anyone ``not of the Caucasian race.'' Black families---the Shelleys in Missouri and the McGhees in Michigan---bought restricted homes from willing sellers, and neighbors sued to oust them.

\section{Procedural History}
The Missouri Supreme Court ordered the Shelleys divested of title; the Michigan courts ordered the McGhees to vacate.

\section{Judicial Votes}
Reversed 6--0. Vinson wrote for Black, Frankfurter, Douglas, Murphy, and Burton; Reed, Jackson, and Rutledge did not participate.

\section{Holding}
Private covenants, voluntarily observed, do not violate the Fourteenth Amendment. But judicial enforcement is state action, and enforcing racial covenants against willing buyers and sellers denies equal protection.

\section{Analysis of Opinions}
\textbf{Majority:} the covenants would be plainly invalid as statutes under \textit{Buchanan v. Warley}. Since \textit{Ex parte Virginia}, courts have been treated as the State, even when they enforce common-law rules through fair proceedings. But for the courts' ``active intervention,'' the families would have lived in their homes. Equal enforcement against white buyers is no defense, because the rights are ``personal rights.'' \textbf{No separate opinions;} \textit{Barrows v. Jackson} later barred damages actions.

\section{Examples: Future Applications}
\textbf{Same-Side:} enforcing a covenant against non-Christian buyers is barred. \textbf{Opposite-Side:} a homeowner's trespass complaint against an unwanted guest enforces her own choice. \textbf{Fence-Sitter:} a will conditioning a bequest on marrying within the faith.

\section{Critique}
Read literally, \textit{Shelley} would make all judicially enforced private choices state action, so courts confine it to race. Wechsler found no neutral principle; scholars offer willing-party coercion or \S~1982 as narrower grounds.

\section{Key Quotations}
\begin{quote}
``That Amendment erects no shield against merely private conduct, however discriminatory or wrongful.''
\end{quote}
\noindent --- Chief Justice Vinson, for the Court, 334 U.S. at 13.

\end{document}
